% !TEX root = ../main.tex
% !TEX program = lualatex
\documentclass[../main.tex]{subfiles}
\IfSubfilesClassLoaded{\externaldocument{\subfix{../build/main}}}{}
% =============================================================================
% CHAPTER 6: PPBE (PLANNING, PROGRAMMING, BUDGETING, AND EXECUTION)
% DESCRIPTION: PPBE purpose, phases, FYDP structure, and Navy PPBE flow.
% =============================================================================

\begin{document}
\IfSubfilesClassLoaded{\chapter{EDO Study Guide}}{}
%====================
% Contents here
%====================

% === PPBE ===
\section{PPBE (Planning, Programming, Budgeting, and Execution)}

\subsection{Role of the PPBE process}
\begin{description}
  \item[\textbf{Why PPBE exists.}] The \ac{dod} uses \ac{ppbe} to identify fiscal needs and allocate resources across the Services and \ac{dod} agencies.
  \item[\textbf{Scale.}] Annual budget exceeds \$885B and supports about 3 million active duty and civilian employees.
  \item[\textbf{Balancing act.}] Weighs competing requirements for limited funds while reconciling internal priorities and external direction from \ac{omb} and Congress.
  \item[\textbf{Output.}] The end product is the \ac{dod} portion of the President's Budget.
\end{description}
~\srcCite{edo-3-1-1-ppbe-2025}

\subsection{PPBE overview}
\begin{description}
  \item[\textbf{Planning.}] Predict the strategic environment and define the capabilities it requires.
  \item[\textbf{Programming.}] Allocate resources to missions and integrate programs to minimize shortfalls and duplication.
  \item[\textbf{Budgeting.}] Produce an annual budget for Congress that links missions to requested funding.
  \item[\textbf{Execution.}] Operate in the real world and make adjustments.
\end{description}
~\srcCite{edo-3-1-1-ppbe-2025}

\subsection{PPBE timeline and the DAS}
\ac{ppbe} is calendar-driven, while the \ac{das} is event-driven. Once a program is initiated, PPBE resource allocations can reshape \ac{das} execution plans. Figure~\ref{fig:ppbe_timeline} shows the PPBE timeline across multiple fiscal years~\autocite{edo-3-1-1-ppbe-2025}.

\begin{figure}[H]
  \centering\color{SteelDeep}
  \includegraphics[width=0.85\linewidth]{PPBE Timeline.png}
  \caption[PPBE timeline across fiscal years]{PPBE timeline across fiscal years. \srcCite{edo-3-1-1-ppbe-2025}}
  \label{fig:ppbe_timeline}
\end{figure}

Figure~\ref{fig:ppbe_platinum_card} shows the Navy PPBE timeframes from the DAU Platinum Card.
\begin{figure}[H]
  \centering\color{SteelDeep}
  \includegraphics[width=0.85\linewidth]{Platinum Card PBBE.png}
  \caption[Timeframes of the Navy PPBE process]{Timeframes of the Navy PPBE process. \srcCite{edo-3-1-1-ppbe-2025}}
  \label{fig:ppbe_platinum_card}
\end{figure}

\subsection{Future Years Defense Program (FYDP)}
\begin{description}
  \item[\textbf{Definition.}] The \ac{dod}'s internal database for tracking resource allocations.
  \item[\textbf{Navy access.}] The Navy uses the Program Budget Information System (\ac{pbis}) for controlled access.
  \item[\textbf{Time horizon.}] Prior Year (PY), Current Year (CY), Budget Year (BY), plus four out-years (POM) and additional force-structure-only years.
  \item[\textbf{Purpose.}] Track \ac{dod} funds, force structure, and personnel strength throughout PPBE.
\end{description}
~\srcCite{edo-3-1-1-ppbe-2025}

\subsection{FYDP structure}
The \ac{fydp} lets users identify, sort, and display \ac{dod} plans by three dimensions: Major Force Program, appropriation, and component/agency~\autocite{edo-3-1-1-ppbe-2025}.

\subsubsection{Major Force Programs}
\begin{description}
  \item[\textbf{Twelve MFPs.}] The FYDP organizes data into 12 \acp{mfp}.
  \item[\textbf{Combat vs support.}] Six \acp{mfp} are combat force programs and six are support programs.
\end{description}
~\srcCite{edo-3-1-1-ppbe-2025}

\subsubsection{Appropriations}
Appropriation categories are the major funding classes by which Congress provides funds for all federal agencies, including the \ac{dod}~\autocite{edo-3-1-1-ppbe-2025}.

\subsubsection{Component or agency}
The \ac{fydp} tracks funding for each military service and other \ac{dod} agencies; each data set may include multiple appropriations and \acp{mfp}~\autocite{edo-3-1-1-ppbe-2025}.

\subsubsection{Program elements}
\begin{description}
  \item[\textbf{PE codes.}] Resources are assigned to an \ac{mfp} using unique \ac{pe} codes.
  \item[\textbf{Granularity.}] \acp{pe} can be as narrow as a single platform or as broad as strategic planning.
  \item[\textbf{Multiple appropriations.}] A single \ac{pe} can include multiple appropriation types (for example, procurement plus \ac{omn} for sustainment).
\end{description}
~\srcCite{edo-3-1-1-ppbe-2025}

\subsection{Planning phase}
\begin{description}
  \item[\textbf{Lead.}] Led by \ac{usdp} with inputs from \ac{cjcs}, \acp{ccmd}, and \ac{osd} staff.
  \item[\textbf{Key products.}] \ac{cpr} and \ac{dpg}.
  \item[\textbf{Focus.}] Threat versus capability gaps, strategy updates, and guidance for programming and budgeting.
\end{description}
~\srcCite{edo-3-1-1-ppbe-2025}

\subsubsection{Planning phase key documents}
\begin{description}
  \item[\textbf{CPR.}] \ac{cjcs} coordinates with the \ac{jroc} and \acp{ccmd} to identify priorities for \ac{dpg} input; signed by \ac{cjcs}.
  \item[\textbf{DPG.}] Drafted by \ac{usdp} and signed by \ac{secdef}; provides force-structure and resource priorities for program development and is the end product of the planning phase.
\end{description}
~\srcCite{edo-3-1-1-ppbe-2025}

\subsection{Programming phase}
\begin{description}
  \item[\textbf{POM development.}] Each department and \ac{dod} agency creates a \ac{pom} that allocates resources across programs over five years starting with the budget year.
  \item[\textbf{OSW program review.}] \ac{osd} conducts a program review; \ac{cape} adjudicates issues and \ac{cjcs} signs the \ac{cpa}.
  \item[\textbf{Decision.}] Deputy \ac{secdef} issues a \ac{pdm} approving POMs as modified by the review.
\end{description}
~\srcCite{edo-3-1-1-ppbe-2025}

\subsection{Navy programming process}
\begin{description}
  \item[\textbf{Top-down fiscal guidance.}] Fiscal controls flow from \ac{omb} to \ac{osd}, \ac{don}, \ac{opnav}, resource sponsors, and then programs.
  \item[\textbf{Bottom-up POM.}] Units propose resources; \ac{n8} integrates sponsor submissions into the Navy tentative POM; \ac{don} integrates Navy and USMC inputs and \ac{secnav} signs the \ac{don} POM.
  \item[\textbf{Issues.}] Proposed changes from previous funding levels are tracked and adjudicated as issues.
\end{description}
~\srcCite{edo-3-1-1-ppbe-2025}

\subsubsection{Resource sponsors}
\begin{longtblr}
  [%
  caption = {Navy resource sponsor responsibilities.},%
  label   = {tab:ppbe-resource-sponsors},%
  entry   = {Navy Resource Sponsors},%
  remark{\tSource} = {\tabCite{edo-3-1-1-ppbe-2025}}
  ]{colspec = {@{} lL @{}}}
  \toprule
  {Sponsor} & {Primary focus} \\
  \midrule
  N1 & Manpower, personnel, training, and education. \\
  N2/N6 & Information warfare. \\
  N3/N5 & Navy warfare development. \\
  N4 & Afloat and ashore readiness. \\
  N8 & Integration of capabilities and resources (leads POM balance). \\
  N9 & Warfare systems. \\
  \bottomrule
\end{longtblr}

\subsubsection{Resource sponsor N9}
\begin{description}
  \item[\textbf{Role.}] N9 advocates the largest portion of Navy \ac{toa} each year, integrating expeditionary, surface, undersea, and air warfare requirements with associated systems, manpower, training, and readiness.
\end{description}
~\srcCite{edo-3-1-1-ppbe-2025}

\begin{longtblr}
  [%
  caption = {N9 warfare systems directorates.},%
  label   = {tab:ppbe-n9-directorates},%
  entry   = {N9 Warfare Systems Directorates},%
  remark{\tSource} = {\tabCite{edo-3-1-1-ppbe-2025}}
  ]{colspec = {@{} lL @{}}}
  \toprule
  {Directorate} & {Focus} \\
  \midrule
  N9I & Warfare Integration Directorate. \\
  N95 & Navy Expeditionary Combat Command. \\
  N96 & Naval Surface Force. \\
  N97 & Naval Submarine Force. \\
  N98 & Naval Air Systems Command. \\
  \bottomrule
\end{longtblr}

\subsection{POM resources}
\begin{description}
  \item[\textbf{Allocation.}] Future resources are allocated to buy programs that deliver planning phase capabilities across the FYDP.
  \item[\textbf{Resource types.}] Quantities, end strength, and money.
\end{description}
~\srcCite{edo-3-1-1-ppbe-2025}

\begin{longtblr}
  [%
  caption = {Example of POM resource types for one fiscal year (F/A-18).},%
  label   = {tab:ppbe-pom-resources},%
  entry   = {POM Resource Example},%
  remark{\tSource} = {\tabCite{edo-3-1-1-ppbe-2025}}
  ]{colspec = {@{} l r r r @{}}}
  \toprule
  {Resource type} & {Quantity} & {Unit cost} & {Subtotal} \\
  \midrule
  Aircraft quantity & 3 & \$75{,}000K each & \$225{,}000K \\
  End strength & 25 & \$120K each & \$3{,}000K \\
  Flight hours & 3{,}000 & \$5K per hour & \$15{,}000K \\
  \midrule
  {Total} & {} & {} & {\$243{,}000K} \\
  \bottomrule
\end{longtblr}

\subsection{Navy POM exercise}
\begin{description}
  \item[\textbf{Baseline and new requirements.}] Start with the current funded program of record, add new strategic requirements, and create offsets to balance the tentative POM.
  \item[\textbf{Integration.}] \ac{n8} integrates resource sponsor inputs into the Navy tentative POM and the \ac{don} POM.
\end{description}
~\srcCite{edo-3-1-1-ppbe-2025}

\subsection{POM example (DDG 1000 combat systems)}
\begin{description}
  \item[\textbf{Stakeholders.}] Resource sponsor OPNAV N96E, ship PM (PMS 500), and combat systems PM (PEO IWS 9.0).
  \item[\textbf{Content.}] Detailed POM slides across four program elements over five fiscal years, covering cybersecurity through spares.
\end{description}
~\srcCite{edo-3-1-1-ppbe-2025}

\subsection{Programming versus budgeting}
\begin{longtblr}[%
    caption = {Programming versus budgeting.},
    label   = {tab:ppbe-prog-vs-budg},
    entry   = {Programming versus Budgeting},
    remark{\tSource} = {\tabCite{edo-3-1-1-ppbe-2025}}%
  ]{colspec = {@{} lLL @{}}}
  \toprule
  {} & {Programming} & {Budgeting} \\
  \midrule
  Purpose & Balanced allocation of limited resources within \ac{dod} & Communicate \ac{don} and \ac{dod} intent outside \ac{dod} and inform the President's Budget \\
  Process & \ac{osd}-developed; differs by Service & Constitutionally derived; mostly common across Services \\
  Timeframe & Budget year plus four out-years (\ac{fydp}) & Budget year only (first year of \ac{fydp}) \\
  \bottomrule
\end{longtblr}

Figure~\ref{fig:ppbe_concurrent_review} illustrates the concurrent program and budget review cycle.
\begin{figure}[H]
  \centering\color{SteelDeep}
  \includegraphics[width=0.85\linewidth]{Platinum Card Concurrent Program Budget Review.png}
  \caption[Concurrent program and budget review]{Concurrent program and budget review. \srcCite{edo-3-1-1-ppbe-2025}}
  \label{fig:ppbe_concurrent_review}
\end{figure}

\subsection{Programming phase key documents}
\begin{description}
  \item[\textbf{POM.}] Time-phased resource allocation to support the \ac{dpg}; reviewed for compliance with \ac{dpg} and the National Military Strategy (NMS); signed by Service Secretaries (\ac{secnav} for the Navy).
  \item[\textbf{PDM.}] \ac{secdef} direction that summarizes decisions and approves the component POM as modified by the program and budget review; issued by the Deputy \ac{secdef}.
\end{description}
~\srcCite{edo-3-1-1-ppbe-2025}

\subsection{Budgeting phase key events and documents}
\begin{description}
  \item[\textbf{Budget Estimate Submission (BES).}] Services submit \ac{bes} or Budget Change Proposals in appropriation format; signed by Service Secretaries.
  \item[\textbf{Budget hearings.}] Joint \ac{osd} and \ac{omb} hearings with advance questions on programs.
  \item[\textbf{Program Budget Decisions (PBDs).}] \acp{pbd} direct budget adjustments; Services prepare reclamas; signed by the Deputy \ac{secdef}.
  \item[\textbf{Major Budget Issues (MBI).}] \ac{mbi} provide the last chance to raise issues to \ac{osd} before the \ac{pb}.
  \item[\textbf{President's Budget (PB).}] \ac{pb} is submitted to Congress (nominally the first Monday in February).
\end{description}
~\srcCite{edo-3-1-1-ppbe-2025}

\subsubsection{PBDs and reclamas}
\begin{description}
  \item[\textbf{PBD origin.}] Issued by USW(FM\&C) analysts after review of the \ac{pom}/\ac{bes} and budget hearing responses.
  \item[\textbf{Reclama.}] Draft \acp{pbd} are vetted for comment; affected parties submit reclamas with operational impacts and business case justification.
  \item[\textbf{Final decision.}] USW(FM\&C) recommends to the Deputy SECWAR, who signs the final \acp{pbd}.
\end{description}
~\srcCite{edo-3-1-1-ppbe-2025}

\subsection{Execution phase}
\begin{description}
  \item[\textbf{Purpose.}] Assess programs of record against actual performance and adjust plans and budgets to minimize near-term financial risk.
  \item[\textbf{Activities.}] Mid-year review, reprogramming, supplemental appropriations, and financial closeout.
\end{description}
~\srcCite{edo-3-1-1-ppbe-2025}

\subsection{Phases of PPBE review}
\begin{description}
  \item[\textbf{Planning.}] Identifies capabilities and defines national defense policies, objectives, and strategies; products are \ac{cpr} and \ac{dpg}.
  \item[\textbf{Programming.}] Allocates resources and translates planning into time-phased requirements; products are \ac{pom} and \ac{pdm}.
  \item[\textbf{Budgeting.}] Reviews proper budget justification and execution; products are \ac{bes}, \acp{pbd}, and the \ac{dod} portion of the \ac{pb}.
  \item[\textbf{Execution.}] Concurrent with program and budget review; develops performance metrics and adjusts resources to meet desired performance.
\end{description}
~\srcCite{edo-3-1-1-ppbe-2025}

\subsection{Summary check}
\begin{description}
  \item[\textbf{Four phases.}] Planning, Programming, Budgeting, Execution.
  \item[\textbf{Planning products.}] \ac{cpr} and \ac{dpg}.
  \item[\textbf{Signatures.}] \ac{cpr} signed by \ac{cjcs}; \ac{dpg} signed by \ac{secdef}; \ac{pom}/\ac{bes} signed by Service Secretaries; \ac{pdm}/\acp{pbd} signed by Deputy \ac{secdef}.
  \item[\textbf{FYDP building blocks.}] \acp{pe} identify functional or organizational entities and their resources.
\end{description}
~\srcCite{edo-3-1-1-ppbe-2025}
%====================
% End of file
%====================
\IfSubfilesClassLoaded{
  \RenewDocumentCommand{\entryname}{}{\textbf{\color{Modern} Acronym}}
  \RenewDocumentCommand{\descriptionname}{}{\textbf{\color{Modern} Definition}}
  \printnoidxglossary[
    type=\acronymtype,
    title=Acronyms,
    style=long-booktabs
  ]}{}
\end{document}
